\documentclass[../main.tex]{subfiles}
\graphicspath{{\subfix{../image/}}}

\begin{document}

\chapter{余切丛}\label{chapter:ch11-余切丛}

在本章中，我们引入一个初等微积分中通常见不到的构造：\textbf{切余向量}，即在点 $p\in M$ 处切空间上的线性泛函。在点 $p$ 处所有余向量构成的空间是一个向量空间，称为 $p$ 处的\textbf{余切空间}；用线性代数的话说，它是 $T_pM$ 的对偶空间。$M$ 上所有点处余切空间的并是一个向量丛，称为\textbf{余切丛}。

切向量给了我们曲线导数的与坐标无关的解释，而流形上实值函数的导数则最自然地解释为切余向量。因此我们把实值函数的\textbf{微分}定义为余向量场（余切丛的光滑截面）；它是梯度的与坐标无关的类似物。然后我们考察余向量场在光滑映射下的行为，并证明光滑映射上域上的余向量场总能拉回到定义域上成为余向量场。

本章后半部分引入余向量场的线积分，它是初等微积分中线积分的自然推广。然后我们考察三类密切相关的余向量场之间的关系：恰当的（即某函数的微分）、保守的（绕闭曲线的线积分为零）、以及闭的（在坐标下满足某微分方程）。这引出微积分基本定理在流形线积分上的深远推广。

\subfile{Chapter-11/section-01.tex}
\subfile{Chapter-11/section-02.tex}
\subfile{Chapter-11/section-03.tex}
\subfile{Chapter-11/section-04.tex}
\subfile{Chapter-11/section-05.tex}
\subfile{Chapter-11/section-06.tex}
\end{document}